\chapter{Grafos III: Avançado}

\section{Componentes Fortemente Conexos (SCC)}

Em um grafo dirigido, um \textbf{componente fortemente conexo} é um subconjunto maximal de vértices onde todo vértice consegue alcançar todo outro vértice no subconjunto.

\textbf{Aplicações:} analisar ciclos de dependência, redes sociais, fluxos em circuitos.

\subsection{Tarjan}

Tarjan encontra todos os SCCs em uma única passada de DFS usando uma pilha.

\textbf{Complexidade:} $O(V + E)$.

\subsection{Quando usar}

\begin{itemize}[nosep]
  \item Encontrar \textbf{ciclos} em grafos dirigidos;
  \item Condensar um grafo em componentes \textbf{fortemente conexos};
  \item Quando precisa de uma \textbf{única passada} eficiente.
\end{itemize}

\begin{exemplodiabox}
Você tem um mapa de cidades com estradas de mão única. Tarjan encontra grupos de cidades onde você pode ir de qualquer uma para qualquer outra (ciclos). Isso ajuda a entender se há deadlocks ou dependências circulares.
\end{exemplodiabox}

\incluircppcaption{algoritmos_avulsos/grafos/tarjan.cpp}{Tarjan}

\subsection{Kosaraju}

Kosaraju também encontra SCCs, usando duas DFS: uma no grafo original, outra no grafo transposto.

\textbf{Complexidade:} $O(V + E)$, mas com constante maior que Tarjan.

\subsection{Quando usar}

\begin{itemize}[nosep]
  \item Quando a implementação de Tarjan parece complexa;
  \item Quando já tem o \textbf{grafo transposto} disponível;
  \item Quando precisa de clareza sobre o algoritmo.
\end{itemize}

\begin{exemplodiabox}
Você faz uma DFS para ver quando cada cidade ``termina'' de ser visitada. Depois, em ordem reversa, você faz outra DFS no mapa de volta (arestas reversas). Cada novo grupo encontrado é um SCC.
\end{exemplodiabox}

\incluircppcaption{algoritmos_avulsos/grafos/kosaraju.cpp}{Kosaraju}

\section{Pontes e Articulações}

Uma \textbf{ponte} é uma aresta cuja remoção desconecta o grafo. Uma \textbf{articulação} é um vértice cuja remoção desconecta o grafo.

\textbf{Aplicações:} identificar pontos críticos em redes, robustez de sistemas.

\subsection{Encontrar pontes}

Usando DFS com valores \texttt{tin} (tempo de descoberta) e \texttt{low} (menor vértice alcançável sem usar aresta de volta).

\textbf{Complexidade:} $O(V + E)$.

\subsection{Quando usar}

\begin{itemize}[nosep]
  \item Encontrar \textbf{pontos críticos} em redes (internet, estradas);
  \item Analisar \textbf{robustez} de um grafo;
  \item Identificar componentes biconexos.
\end{itemize}

\begin{exemplodiabox}
Você tem um mapa de cidades conectadas por estradas. Uma ponte é uma estrada cuja destruição deixaria cidades isoladas. Encontrar pontes ajuda a reforçar a infraestrutura.
\end{exemplodiabox}

\incluircppcaption{algoritmos_avulsos/grafos/pontes.cpp}{Encontrar pontes}

\section{Fluxo Máximo}

O \textbf{fluxo máximo} entre dois vértices é o máximo de ``carga'' que consegue passar de um ponto a outro respeitando as capacidades das arestas.

\textbf{Teorema de Max-Flow Min-Cut:} o fluxo máximo é igual à capacidade mínima do corte que separa origem e destino.

\textbf{Aplicações:} roteamento de rede, matching em grafos bipartidos, circulação de recursos.

\subsection{Edmonds-Karp}

Edmonds-Karp é Ford-Fulkerson com BFS para encontrar caminhos de aumento.

\textbf{Complexidade:} $O(VE^2)$.

\subsection{Quando usar}

\begin{itemize}[nosep]
  \item Quando os números não são muito grandes;
  \item Quando precisa de uma implementação \textbf{simples e correta};
  \item Para verificação rápida em competições.
\end{itemize}

\begin{exemplodiabox}
Você tem um sistema de canos com capacidades diferentes. O fluxo máximo é quanto você consegue bombear de ponta a ponta. Edmonds-Karp encontra isso achando caminhos de aumento iterativamente.
\end{exemplodiabox}

\incluircppcaption{algoritmos_avulsos/grafos/edmonds_karp.cpp}{Edmonds-Karp}

\subsection{Dinic}

Dinic é uma versão otimizada de Ford-Fulkerson que usa grafos de níveis.

\textbf{Complexidade:} $O(V^2 E)$, muito mais rápido na prática.

\subsection{Quando usar}

\begin{itemize}[nosep]
  \item Quando os números são \textbf{muito grandes};
  \item Quando capacidades são altas (Edmonds-Karp fica lento);
  \item Na maioria dos problemas de \textbf{fluxo em competições}.
\end{itemize}

\begin{exemplodiabox}
Dinic é Edmonds-Karp turbinado: em vez de achar um caminho por vez, ele encontra múltiplos caminhos em paralelo usando um grafo auxiliar de níveis.
\end{exemplodiabox}

\incluircppcaption{algoritmos_avulsos/grafos/dinic.cpp}{Dinic}

\section{Min-Cut}

O \textbf{corte mínimo} é o conjunto de arestas cuja remoção desconecta origem do destino, com capacidade total mínima.

\textbf{Teorema:} Min-Cut = Max-Flow (em grafos com capacidades).

\subsection{Quando usar}

\begin{itemize}[nosep]
  \item Encontrar o \textbf{gargalo crítico} de uma rede;
  \item Particionar um grafo com \textbf{mínimo de perda};
  \item Aplicações em visão computacional (segmentação de imagens).
\end{itemize}

\begin{exemplodiabox}
Se você quer cortar uma rede de canos em dois grupos com perda mínima de capacidade, o corte mínimo identifica exatamente quais canos cortar. É tão importante que tem nome matemático próprio.
\end{exemplodiabox}

\incluircppcaption{algoritmos_avulsos/grafos/mincut.cpp}{Min-Cut}

\section{2-SAT}

\textbf{2-SAT} é um problema de satisfazibilidade booleana onde cada cláusula tem exatamente 2 literais. \textbf{Solúvel em tempo polinomial} usando grafos de implicação e SCCs.

\textbf{Complexidade:} $O(V + E)$.

\subsection{Como funciona}

Cada variável $x_i$ vira dois vértices: $x_i$ e $\neg x_i$ (negação). Cada cláusula $(a \vee b)$ vira implicações: $\neg a \Rightarrow b$ e $\neg b \Rightarrow a$. Depois, encontra SCCs: se $x_i$ e $\neg x_i$ estão no mesmo SCC, o problema é \textbf{insatisfazível}.

\subsection{Quando usar}

\begin{itemize}[nosep]
  \item Problemas com \textbf{restrições booleanas} (sim/não);
  \item Agendamento com conflitos;
  \item Coloração de grafos com 2 cores e restrições.
\end{itemize}

\begin{exemplodiabox}
Você tem 3 tarefas e cada uma pode ser feita de dia ou de noite, com restrições tipo ``Se A é de dia, então B deve ser de noite''. 2-SAT descobre se existe uma atribuição válida e retorna ela.
\end{exemplodiabox}

\incluircppcaption{algoritmos_avulsos/grafos/two_sat.cpp}{2-SAT}
